\documentclass[10pt,a4paper]{amsart}
\usepackage[margin=19mm]{geometry}
\usepackage{amsmath,amssymb,mathtools}
\usepackage[scr=rsfs,cal=euler]{mathalfa}
\usepackage{xfrac}
\usepackage[unicode,hidelinks,bookmarks=false]{hyperref}
\newcommand{\ie}{i.e.\ }
\newcommand{\cf}{cf.\ }
\newcommand{\resp}{respectively\ }
\newcommand{\Subsection}{\subsection}
\providecommand{\nobreakdash}{\nobreak\hbox{-}}
\setlength{\emergencystretch}{2em}
\multlinegap0pt
\usepackage{bm}
\usepackage{bbm}
\usepackage{dsfont}
\usepackage{xspace}
\usepackage{cancel}
\usepackage{mathtools}
\usepackage{amsthm}
\makeatletter
\@ifundefined{theorem}{\newtheorem{theorem}{Theorem}}{}
\@ifundefined{prop}{\newtheorem{prop}[theorem]{Proposition}}{}
\makeatother
\makeatletter
\expandafter\ifx\csname myeq\endcsname\relax
\newenvironment{myeq}[1][]
{\stepcounter{theorem}\begin{equation}\tag{\thetheorem}{#1}}
{\end{equation}}
\fi
\expandafter\ifx\csname myeqa\endcsname\relax
\newenvironment{myeqa}[1][]
%{\stepcounter{theorem}\begin{equation}\tag{\thetheorem}{#1}\begin{aligned}}
%{\end{aligned}}


\fi
\expandafter\ifx\csname mysubsection\endcsname\relax
\newenvironment{mysubsection}[2][]
{\begin{subsec}\begin{upshape}\begin{bfseries}{#2.}
\end{bfseries}{#1}}
{\end{upshape}\end{subsec}}
\fi
\renewcommand{\rho}{\varrho}
\newcommand{\uZ}{\underline{\mathbb{Z}}}
\makeatother
\title[Revisiting the Nandakumar-Ramana Rao Conjecture]{Revisiting the Nandakumar-Ramana Rao Conjecture}
\author{Surojit Ghosh, Ankit Kumar}
\date{Source: arXiv:2510.25186v1 (CC BY 4.0)}
\begin{document}
\maketitle

\noindent\emph{Source and licence.} Revisiting the Nandakumar-Ramana Rao Conjecture. Version arXiv:2510.25186v1 (CC BY 4.0), distributed under the Creative Commons Attribution 4.0 International licence, as recorded in the arXiv licence metadata for this exact version and shown on the abstract page https://arxiv.org/abs/2510.25186v1. Full rights notice: licenses/arxiv-2510.25186-CC-BY-4.0.txt.\par\smallskip

\noindent\textit{Editorial scope.} Counted unit: the Prop whose source label is \texttt{None} in the article (source0.tex lines 509--511, source label \texttt{None}), delivered together with the article's own complete original proof of it (source0.tex lines 513--535, one \texttt{proof} environment of 23 lines). No mathematics is written, repaired or re-derived: statement and proof are the article's own text line for line. Required context retained uncounted: none. Every definition, display and label of those lines is retained unchanged. Omitted from this package and not redistributed: 1--508, 512--512, 536--642. Nothing inside a retained excerpt is omitted or altered: every retained body is delivered exactly as the article's lines are written, so no omission receipt of any kind accompanies this package. Citations: the retained excerpts carry no citation command at all, so no bibliography entry of the article is transcribed and no reference is redistributed. Reference adaptation: none was needed and none was made; every reference inside the delivered unit resolves inside it, so no unresolved reference or citation is produced. Printed slips kept literal: no printed word, symbol, hypothesis or number of the counted statement or of its proof was altered. Layout and class: the article's own class is \texttt{amsart} with packages including amsfonts, amssymb, verbatim, amsmath, amsthm, latexsym, textcomp, amscd, epsfig, graphics, graphicx, color, url, enumerate; the wrapper is the lane's pinned \texttt{amsart} editorial wrapper at 10pt on A4, which restates the theorem-like environments the retained text needs and adds the article's own macro definitions that the retained text needs (\texttt{\textbackslash rho}, \texttt{\textbackslash uZ}), with extra wrapper packages (bm, bbm, dsfont, xspace, cancel, mathtools). The delivered numbering is the unit's own. Assets: no retained span contains a figure environment, a graphics-inclusion command, a PDF-page inclusion, a table or a TikZ picture; no image file is copied into this package, the package declares no asset dependency and no image file is redistributed with it. Licence: version arXiv:2510.25186v1 of this article is distributed under the Creative Commons Attribution 4.0 International licence, as recorded in the arXiv licence metadata for this exact version and shown on the abstract page https://arxiv.org/abs/2510.25186v1. A full rights notice is delivered at licenses/arxiv-2510.25186-CC-BY-4.0.txt. Characters: every retained excerpt is pure ASCII, so the delivered engine, XeLaTeX with its default Latin Modern text font, reports no missing character.
\begin{prop}
    For a cyclic group $G$ with order having two distinct prime divisors, the Euler class $a_{\bar{\rho}}$ vanishes. Here, $\bar{\rho}$ denotes the reduced regular representation of $G.$
\end{prop}

\begin{proof}
Let \( G = C_n \) be the cyclic group of order \( n \), and let \( p \neq q \) be two distinct prime divisors of \( n \). We construct two subrepresentations \( V_p \) and \( V_q \) of the reduced regular representation \( \bar{\rho} \) such that
\[
\gcd\left( \frac{|G|}{|G_{V_p}|}, \frac{|G|}{|G_{V_q}|} \right) = 1,
\]
where \( G_V \) denotes the stabilizer subgroup of the representation \( V \).

To achieve this, let \( \eta \) be a non-trivial one-dimensional complex \( G \)-representation induced by a group homomorphism \( \eta \colon G \to S^1 \), with kernel \( H = \ker(\eta) \). The corresponding representation sphere \( S^\eta \) admits a \( G \)-CW complex structure of the form
\[
S^0 \cup (G/H_+ \wedge e^1) \cup (G/H_+ \wedge e^2).
\]
The associated equivariant cellular chain complex with constant coefficients \( \uZ \) is
\[
\mathbb{Z} \xrightarrow{0} \mathbb{Z} \xrightarrow{|G/H|} \mathbb{Z},
\]
which gives rise to the isomorphism
\[
\tilde{H}^\eta_G(S^0; \uZ) \cong \mathbb{Z}/|G/H|\mathbb{Z}.
\]
In particular, this implies that \( |G/H| \cdot a_\eta = 0 \).

We now choose the subrepresentations \( V_p = \xi^{n/p} \) and \( V_q = \xi^{n/q} \), where \( \xi^d \) denotes the one-dimensional complex representation of \( G \) corresponding to rotation by angle \( \frac{2\pi d}{n} \). Since the indices \( \frac{|G|}{|G_{V_p}|} \) and \( \frac{|G|}{|G_{V_q}|} \) are coprime, it follows that the Euler class \( a_{\bar{\rho}} = 0 \).
\end{proof}

\end{document}
